% ================================================================================
% Section 6: Results
% ================================================================================
\section{Results}

\subsection{Foundational Results}

\begin{frame}{Proposition 1 - Key Insight for Computability}
    \begin{proposition}
       Equation $1+p^a = q^b + r^c$ is computable if at least one of these conditions holds:
       \begin{enumerate}
           \item There exists $M$ such that $a \leq M$ for any solution $(a,b,c)$
           \item There exist $M_b, M_c$ such that $b \leq M_b$ and $c \leq M_c$ for any solution
       \end{enumerate}
    \end{proposition}

    \textbf{Proof sketch (Condition 1):}
    \begin{align*}
        q^b+r^c &= 1+p^a \\
        q^b &= 1+p^a - r^c \leq p^a \leq p^M \\
        \implies b &\leq \frac{M\ln p}{\ln q},\quad c \leq \frac{M\ln p}{\ln r}
    \end{align*}

    \textbf{Conclusion:} All solutions lie in a finite set, hence computable in finite time.
\end{frame}

\begin{frame}{Modeling the Modulo Trick with Sets}
    \begin{block}{What happens when we apply the modulo trick?}
        \begin{itemize}
            \item We start with \textbf{all possible} exponent triples $(a,b,c)$
            \item We choose a modulus $m_1$ and impose: $1+p^a \equiv q^b+r^c \pmod{m_1}$
            \item This eliminates triples that don't satisfy this congruence
            \item We repeat with $m_2, m_3, ...$ - each time further restricting the possibilities
        \end{itemize}
    \end{block}

    \vspace{0.3cm}
    \textbf{We can represent this mathematically:}
    \begin{align*}
        S_0^{p,q,r} &= \mathbb{Z}_{\geq 0}^3 \quad \text{(all possibilities)} \\
        S_1^{p,q,r} &= \{(a,b,c) \in S_0^{p,q,r} \; | \; 1+p^a \equiv q^b+r^c \pmod{m_1}\} \\
        S_2^{p,q,r} &= \{(a,b,c) \in S_1^{p,q,r} \; | \; 1+p^a \equiv q^b+r^c \pmod{m_2}\} \\
        &\vdots \\
        S_k^{p,q,r} &= \{(a,b,c) \in S_{k-1}^{p,q,r} \; | \; 1+p^a \equiv q^b+r^c \pmod{m_k}\}
    \end{align*}
\end{frame}

\begin{frame}{From Sets to Definition of Solvability}
    \begin{block}{Key observation}
        After applying all moduli $m_1,...,m_k$, the set $S_k^{p,q,r}$ contains \textbf{all potential solutions} that survive all congruence conditions.
    \end{block}

    \vspace{0.2cm}
    \textbf{Recall Proposition 1:} If we can bound $a$ (or bound both $b$ and $c$), then all variables are bounded and the equation becomes computable by brute force.

    \vspace{0.3cm}
    \textbf{Therefore, the modulo trick succeeds exactly when:}
    \begin{itemize}
        \item After some finite sequence of moduli, the set $S_k^{p,q,r}$ contains only triples where:
        \begin{enumerate}
            \item $a$ is bounded by some $M_a$, OR
            \item both $b$ and $c$ are bounded by some $M_b, M_c$ respectively
        \end{enumerate}
    \end{itemize}

    \vspace{0.2cm}
    This is precisely our definition of \textbf{solvable with modulo trick} - it directly captures what the method aims to achieve!
\end{frame}

\begin{frame}{Formal Definition of Modulo Trick Solvability}
    Let $m_1,m_2,...,m_k$ be the sequence of moduli used. Define:
    \[S_0^{p,q,r} = \mathbb{Z}_{\geq 0}^3\]
    \[S_i^{p,q,r} = \{(a,b,c)\in S_{i-1}^{p,q,r} \; | \; 1+p^a \equiv q^b+r^c \pmod{m_i}\}\]

    \begin{Definitionn}
        Equation \eqref{general} is \textbf{solvable with modulo trick} if $\exists$ finite moduli $m_1,...,m_k$ such that:
        \begin{itemize}
            \item $\exists M_a$ with $M_a \geq a$ for all $(a,b,c) \in S_k^{p,q,r}$, OR
            \item $\exists M_b,M_c$ with $M_b \geq b$ and $M_c \geq c$ for all $(a,b,c) \in S_k^{p,q,r}$
        \end{itemize}
    \end{Definitionn}

    \begin{Definitionn}
        Equation \eqref{general} is \textbf{non-solvable with modulo trick} if for any finite moduli $m_1,...,m_k$:
        \begin{itemize}
            \item $\forall M_a$, $\exists (a,b,c) \in S_k^{p,q,r}$ with $a > M_a$
            \item $\forall M_b,M_c$, $\exists (a,b,c) \in S_k^{p,q,r}$ with $b > M_b$ or $c > M_c$
        \end{itemize}
    \end{Definitionn}
\end{frame}

\begin{frame}{Key Lemma - Single modulus lemma}
    \begin{lemmas}
        Let $A_m = \{(a,b,c) \in \mathbb{Z}_{\geq 0}^3 \; | \; 1+p^a \equiv q^b+r^c \pmod{m}\}$.
        For any moduli $m_1,...,m_k$,
        \[S_k^{p,q,r} = \bigcap_{i=1}^k A_{m_i} = A_{\operatorname{lcm}(m_1,...,m_k)}\]
    \end{lemmas}

    \textbf{Significance:}
    \begin{itemize}
        \item Any sequence of modular restrictions reduces to a \textbf{single modulus} (the LCM)
        \item Simplifies analysis dramatically
        \item Allows us to prove non-solvability by considering arbitrary single modulus
    \end{itemize}
\end{frame}

\subsection{Main Theorems}

\begin{frame}{Theorem 1 - Non-solvability of $1+(pq)^a = p^b + q^c$}
    \begin{theorems}
        For fixed coprime positive integers $p,q$, the equation
        \[1+(pq)^a = p^b + q^c\]
        is non-solvable with modulo trick.
    \end{theorems}

    \textbf{Proof strategy:}
    \begin{enumerate}
        \item By Lemma 3, it suffices to show for \textbf{any} modulus $m$, $A_m$ contains arbitrarily large $a,b,c$
        \item Factor $m = \chi \cdot \Gamma \cdot \mathrm{Z}$ where:
            \begin{itemize}
                \item $\chi$ contains primes dividing $p$
                \item $\Gamma$ contains primes dividing $q$
                \item $\mathrm{Z}$ contains primes dividing neither
            \end{itemize}
        \item Construct solutions using Euler's theorem
        \item Parameters can be arbitrarily large $\implies$ non-solvable
    \end{enumerate}
\end{frame}

\begin{frame}{Corollary 1 - Validating Brenner and Foster (8.037)}
    \begin{corollarys}
        For fixed distinct primes $p,q$, the equation
        \[1+(pq)^a = p^b + q^c\]
        is non-solvable with modulo trick.
    \end{corollarys}

    \vspace{0.3cm}
    \textbf{Significance:}
    \begin{itemize}
        \item This is exactly the case mentioned by Brenner and Foster
        \item They observed it "seems not amenable" to congruence methods
        \item We have now \textbf{proved} this observation rigorously
        \item Provides mathematical validation for their heuristic comment
    \end{itemize}
\end{frame}

\subsection{Generalization}

\begin{frame}{Moving Beyond the Specific Form}
    \begin{block}{So far we've focused on:}
        \[1+p^a = q^b + r^c\]
    \end{block}

    \begin{itemize}
        \item This is just ONE type of exponential Diophantine equation
        \item We want to analyze other forms like:
        \begin{align*}
            &p^x - p^y = q^z - q^w \quad \text{(Brenner \& Foster 8.033)}\\
            &(pq)^a - p^b - q^c + 1 = 0 \quad \text{(Brenner \& Foster 8.037)}
        \end{align*}

        \item \textbf{Significance:} To prove non-solvability for these cases, we need definitions that work for ANY exponential Diophantine equation

        \item \textbf{Generalization strategy:} Extend our framework to handle equations of the form:
        \[C_1 + \sum_{i=1}^n k_i a_i^{x_i} = C_2 + \sum_{j=1}^m l_j b_j^{y_j}\]
    \end{itemize}
\end{frame}

\begin{frame}{General Form of Exponential Diophantine Equation}
    \begin{definition}[General Form]
        Let $n,m \in \mathbb{Z}^{+}$, and $C_1, C_2 \in \mathbb{Z}_{\geq 0}$ and $k_i, a_i, l_j, b_j \in \mathbb{Z}^{+}$ for $i=1,2,..,n$ and $j=1,2,...,m$.

        Consider the exponential Diophantine equation
        \begin{equation}
            C_1 + \sum_{i=1}^n k_i a_i^{x_i} = C_2 + \sum_{j=1}^m l_j b_j^{y_j} \label{general}
        \end{equation}
        where $x_1,x_2,..,x_n, y_1,y_2,..,y_m \in \mathbb{Z}_{\geq 0}$.
    \end{definition}

    \textbf{Examples rewritten in this form:}
    \begin{itemize}
        \item $1 + 5^a = 7^b + 7^c$
        \item $p^x - p^y = q^z - q^w$
    \end{itemize}

    \textbf{Significance:} This notation lets us analyze different exponential Diophantine equations with one common framework.
\end{frame}

\begin{frame}{Generalized Computability Criterion}
    \begin{lemma}[General Computability]
        Equation \eqref{general} is computable if at least one of these conditions hold:
        \begin{enumerate}
            \item There exist positive integer $M$ such that $x_1, x_2, ..., x_n \leq M$ for any solution.
            \item There exist positive integer $M$ such that $y_1, y_2, ..., y_m \leq M$ for any solution.
        \end{enumerate}
    \end{lemma}

    \textbf{Significance of this lemma:}
    \begin{itemize}
        \item This generalizes Proposition 1 from our specific equation to ANY equation of form \eqref{general}
    \end{itemize}
\end{frame}

\begin{frame}{Generalized Solution Sets After Moduli}
    \begin{definition}[Solution Sets]
        Let $m_1, m_2, ..., m_k$ be the sequence of moduli used for the modulo trick. We define:

        \[S_0 = \mathbb{Z}_{\geq 0}^{n+m}\]

        and for each $j = 1,2,...,k$:

        \[S_j = \left\{(x_1, x_2,..., x_n, y_1, y_2,..., y_m) \in S_{j-1}
        \;\middle|\; C_1 + \sum_{i=1}^n k_i a_i^{x_i} \equiv C_2 + \sum_{i=1}^m l_i b_i^{y_i} \pmod{m_j} \right\}\]
    \end{definition}

    \textbf{Notation:} For any modulus $k$, define the congruence solution set:

    \[A'_k = \left\{(x_1,...,x_n,y_1,...,y_m) \in \mathbb{Z}_{\geq 0}^{n+m}  \; |  \; C_1 + \sum_{i=1}^n k_i a_i^{x_i} \equiv C_2 + \sum_{i=1}^m l_i b_i^{y_i} \pmod{k} \right\}\]

    \textbf{Significance:} $S_k$ represents all tuples that survive all congruence conditions after applying the modulo trick. Any true solution to the original equation must be contained in $S_k$.
\end{frame}

\begin{frame}{General Definition of Solvability with Modulo Trick}
    \begin{definition}[General Solvability]
        The exponential Diophantine equation of the form \eqref{general} is called \textbf{solvable with modulo trick} if there exists a finite sequence of moduli $m_1, m_2, ..., m_k \in \mathbb{Z}^{+}$ used for modulo trick such that one of these conditions holds:

        \begin{itemize}
            \item There exists an integer $M$ such that $M \geq x_1, x_2, ..., x_n$ for any $(x_1, x_2, ..., x_n, y_1, y_2, ..., y_m) \in S_k$.

            \item There exists an integer $M$ such that $M \geq y_1, y_2, ..., y_m$ for any $(x_1, x_2, ..., x_n, y_1, y_2, ..., y_m) \in S_k$.
        \end{itemize}
    \end{definition}

    \textbf{Significance of this definition:}
    \begin{itemize}
        \item Captures exactly what the modulo trick aims to achieve - bounding one side of the equation
        \item Therefore, only finitely many tuples remain to check - the equation becomes computable!
        \item This definition works for ANY equation of form \eqref{general}
    \end{itemize}
\end{frame}

\begin{frame}{General Definition of Non-solvability with Modulo Trick}
    \begin{definition}[General Non-solvability]
        The exponential Diophantine equation of the form \eqref{general} is called \textbf{non-solvable with modulo trick} if for each finite sequence of moduli $m_1, m_2, ..., m_k \in \mathbb{Z}^{+}$ used for modulo trick, the following conditions hold for every $N, M \in \mathbb{Z}^{+}$:

        \begin{itemize}
            \item There exists $(x_1, x_2, ..., x_n, y_1, y_2, ..., y_m) \in S_k$ such that $N < x_i$ for some $i = 1,2,...,n$.

            \item There exists $(x_1, x_2, ..., x_n, y_1, y_2, ..., y_m) \in S_k$ such that $M < y_i$ for some $i = 1,2,...,m$.
        \end{itemize}
    \end{definition}

    \textbf{Significance of this definition:}
    \begin{itemize}
        \item No matter which moduli we choose, we cannot force either side of the equation to be bounded
        \item This formalizes what Brenner and Foster meant by "not amenable" to their method
    \end{itemize}
\end{frame}

\begin{frame}{Key Lemma - Generalized Single modulus lemma}
    \begin{lemmas}
        For any moduli $m_1,...,m_k$,
        \[S_k = \bigcap_{i=1}^k A'_{m_i} = A'_{\operatorname{lcm}(m_1,...,m_k)}\]
    \end{lemmas}

    \textbf{Significance:}
    \begin{itemize}
        \item Generalized of lemma 1
        \item Any sequence of modular restrictions reduces to a \textbf{single modulus} (the LCM)
        \item Allows us to prove non-solvability by considering arbitrary single modulus
    \end{itemize}
\end{frame}

\begin{frame}{Lemma 3 - Period of Exponent Residues}
    \begin{lemmas}
        Let $a,n$ be positive integers. The sequence $(a^x \bmod n)_{x=1}^\infty$ is eventually periodic.

        That is, $\exists$ non-negative integers $k,\epsilon$ such that
        \[a^x \equiv a^{x+\epsilon} \pmod n\]
        for all $x \geq k$.
    \end{lemmas}

    \textbf{Proof sketch:}
    \begin{itemize}
        \item Residues are in $\{0,1,...,n-1\}$ (finite set)
        \item By pigeonhole principle, some residue repeats: $a^i \equiv a^{i+\epsilon}$
        \item Multiply by $a^t$ to extend periodicity
        \item \textbf{Key insight:} Exponents can grow while maintaining congruence patterns
    \end{itemize}
\end{frame}

\begin{frame}{Theorem 2 - Non-solvability of $a^x-a^y = b^z-b^w$}
    \begin{theorems}
        Let $a,b \in \mathbb{Z}^+$. The exponential Diophantine equation
        \[a^x - a^y = b^z - b^w\]
        is non-solvable with modulo trick.
    \end{theorems}

    \textbf{Proof sketch:}
    \begin{enumerate}
        \item For any modulus $d$, by Lemma 3, $\exists$ periods $\epsilon_a,\epsilon_b$ and starting points $k_a,k_b$
        \item For any $x \geq k_a$, $y \geq k_b$,
            \[(x,\; x+N\epsilon_a,\; y,\; y+M\epsilon_b) \in A'_d\]
        \item$x+N\epsilon_a,\;y+M\epsilon_b $ are increasing for any positive integer $N,M$
        \item Thus arbitrarily large exponents exist in $A'_d$ for any $d$
        \item Therefore non-solvable by modulo trick
    \end{enumerate}
\end{frame}

\begin{frame}{Corollary 2 - Validating Brenner and Foster (8.033)}
    \begin{corollarys}
        For fixed distinct primes $p,q$, the equation
        \[p^x - p^y = q^z - q^w\]
        is non-solvable with modulo trick.
    \end{corollarys}

    \vspace{0.3cm}
    \textbf{Significance:}
    \begin{itemize}
        \item Validates Brenner and Foster's comment (8.033)
        \item They noted this equation "cannot be fully investigated by reducing it to congruences"
        \item Our proof shows why: residues become periodic, allowing arbitrarily large solutions to satisfy any congruence
    \end{itemize}
\end{frame}
